% !TeX root = ../main.tex
\chapter{Modelo de ruido y desempe\~no te\'orico del receptor}
\label{ap:ruido}

Este ap\'endice re\'une el modelo de ruido del receptor preamplificado, el
requisito de OSNR que impone el umbral HD-FEC. El \cref{cap:enlace} toma estos resultados como datos para el balance de potencia.


\section{Modelo de ruido y OSNR}
\label{sec:enl-ruido}

La fuente de ruido dominante del enlace es la emisi\'on espont\'anea
amplificada del EDFA del receptor. El modelo est\'a
desarrollado en la \cref{sec:canal-ase}: densidad espectral por modo de
polarizaci\'on $S_{\mathrm{pol}} = \nsp h\nu (G-1) = \tfrac12 (FG-1)h\nu$,
inyectada en cada polarizaci\'on, y OSNR definida sobre el total en
$\Bref = \SI{12.5}{\giga\hertz}$.

\begin{puntos}
  \item Consistencia del factor de emisi\'on espont\'anea:
En la simulaci\'on no se toma $\nsp$ como par\'ametro sino que se usa directamente la f\'ormula
$S_{\mathrm{pol}} = \tfrac12(FG-1)h\nu$. Invirtiendo la
definici\'on de figura de ruido,
\begin{equation}
  F = \frac{1 + 2\nsp(G-1)}{G}
  \quad\Longrightarrow\quad
  \nsp = \frac{FG - 1}{2\,(G-1)},
  \label{eq:nsp-exacto}
\end{equation}
y con $G = \SI{25}{\dB}$ y $F = \SI{4.5}{\dB}$ resulta $\nsp = \num{1.4121}$.
Utilizando la aproximaci\'on de ganancia alta se obtiene: $\nsp \simeq F/2 = \num{1.4092}$. Es decir,
a \SI{25}{\dB} de ganancia las dos expresiones son equivalentes.

El factor de emisi\'on espont\'anea
es $\nsp = N_2/(N_2 - N_1)$, con $N_2$ y $N_1$ las poblaciones de los niveles
superior e inferior del medio activo, de modo que siempre se cumple $\nsp \ge 1$. A su vez, la igualdad
corresponde a inversi\'on de poblaci\'on completa, resultante en
$F = \SI{3.01}{\dB}$, es decir, el l\'imite cu\'antico de un amplificador \'optico. Los
\SI{4.5}{\dB} especificados est\'an entonces
\begin{equation}
  F - \SI{3.01}{\dB} = 10\log_{10}\nsp = \SI{1.50}{\dB}
  \label{eq:exceso-nf}
\end{equation}
por encima del l\'imite, que es lo t\'ipico de un EDFA comercial.
El modelo de
ruido no tiene par\'ametros libres. Fijados $G$ y $F$, la densidad de ASE queda
determinada. 

  \item En la simulaci\'on solo se tiene en cuenta el ruido ASE del EDFA del receptor,
 despreciando al ruido shot, t\'ermico y ASE del booster del transmisor.

En un receptor coherente todos los t\'erminos de ruido aparecen en el
fotodetector multiplicados por el oscilador local, as\'i que la comparaci\'on
correcta es entre densidades espectrales de corriente a la salida de un
fotodiodo. Con $R$ la responsividad, $P_{LO}$ la potencia de oscilador local
que llega a ese fotodiodo, $S_{\mathrm{pol}}$ la densidad de ASE
co-polarizado y $F_{\mathrm{TIA}}$ la figura de ruido el\'ectrica del amplificador de
transimpedancia, el batido ASE--LO, el ruido de shot y el ruido t\'ermico
valen
\begin{equation}
  S_{\mathrm{ASE\text{-}LO}} = 2R^{2}P_{LO}S_{\mathrm{pol}},
  \qquad
  S_{\mathrm{shot}} = 2qRP_{LO},
  \qquad
  S_{\mathrm{term}} = \frac{4k_BT F_{\mathrm{TIA}}}{R_L}.
  \label{eq:ruidos}
\end{equation}

\begin{comment}
Conviene no confundir $F_{\mathrm{TIA}}$ con la $F$ del EDFA. La del EDFA es una figura de
ruido \'optica, mide cu\'anto degrada el amplificador la relaci\'on se\~nal a
ruido del campo \'optico y tiene un piso cu\'antico de \SI{3.01}{\dB}. $F_{\mathrm{TIA}}$
es una figura de ruido el\'ectrica: el factor por el que el amplificador de
transimpedancia multiplica al ruido t\'ermico de la resistencia de carga, sin
m\'as piso que la unidad, y vale t\'ipicamente entre \SI{2}{\dB} y
\SI{4}{\dB} en un TIA de banda ancha.
\end{comment}

El cociente entre los dos primeros t\'erminos no depende de la potencia del
oscilador local:
\begin{equation}
  \frac{S_{\mathrm{shot}}}{S_{\mathrm{ASE\text{-}LO}}}
    = \frac{q}{R\,S_{\mathrm{pol}}}
    = \num{2.81e-3} \equiv \SI{-25.5}{\dB}.
  \label{eq:shot-vs-ase}
\end{equation}
Aumentar la ganancia aumenta $S_{\mathrm{pol}}$ y vuelve despreciable al ruido
shot. Para calcular la relaci\'on con el ruido t\'ermico hay que
definir valores. Con $P_{LO} = \SI{10}{\dBm}$ repartidos entre dos
polarizaciones y las cuatro salidas del h\'ibrido, se tienen \SI{1.25}{\milli\watt} por
fotodiodo, y con $R_L = \SI{50}{\ohm}$, $T = \SI{300}{\kelvin}$ y
$F_{\mathrm{TIA}} = \SI{3}{\dB}$ su densidad espectral termina \SI{23.3}{\dB} por debajo del batido
ASE--LO. Sumando los ruidos t\'ermico y shot, el ruido total sube un factor \num{1.0075}, es decir
\SI{0.03}{\dB}.


Tambi\'en es despreciado el ASE del booster de transmisi\'on. Este ruido se genera
antes del enlace, de modo que llega al receptor atenuado por los
\SI{73.40}{\dB} de p\'erdida de trayecto, y el pre-amplificador lo vuelve a
amplificar. Sin embargo, antes de descartarlo es necesario verificar si este ruido re-amplificado es comparable con el ASE del preamplificador. Con $G_b$ y $F_b$ la ganancia y
la figura de ruido del booster, y $L$ la p\'erdida del enlace,

\begin{equation}
  \frac{S_{\mathrm{booster}}}{S_{\mathrm{pol}}}
    = \frac{\tfrac12\,(F_bG_b-1)\,h\nu\;L\;G}{\tfrac12\,(FG-1)\,h\nu}
    \;\simeq\; \frac{F_b G_b}{F}\,L .
  \label{eq:ase-booster}
\end{equation}

Con $G_b = \SI{35}{\dB}$ y $F_b = \SI{5}{\dB}$, el
ASE del booster comienza $F_bG_b/F = \SI{35.5}{\dB}$ por encima de la densidad del EDFA del receptor,
porque el booster tiene \SI{10}{\dB} m\'as de ganancia. Sin embargo, luego de pasar por el canal y amplificarse, queda \SI{37.9}{\dB} por
debajo, o \SI{0.001}{\dB} de contribuci\'on.


Todo esto vale para el receptor preamplificado del enlace. El
$S_{\mathrm{pol}}$ de la \cref{eq:ruidos} crece con la ganancia del
pre-amplificador. Por lo tanto el propio EDFA termina generando el ruido que hace despreciable a los dem\'as. Alcanza con \SI{2.6}{\dB} de ganancia para que
el batido ASE--LO supere al t\'ermico, y a \SI{25}{\dB} ya le saca
\SI{23}{\dB}. En el esquema experimental del \cref{cap:experimental} no se tiene el
pre-amplificador, por lo que los ruidos dominantes son tanto el shot como el t\'ermico, dependiendo
de la potencia del LO.

 La \cref{tab:ruidos} resume los resultados.
\end{puntos}

\begin{table}[htbp]
  \centering

  \small
  \setlength{\tabcolsep}{4pt}
  \begin{tabular}{lccl}
    \toprule
    Fuente & Densidad relativa & Penalidad & En el modelo \\
    \midrule
    Batido se\~nal--ASE del pre-amplificador & referencia      & ---                & s\'i \\
    Batido ASE--ASE                          & ---             & ---                & s\'i \\
    Ruido t\'ermico del TIA                  & \SI{-23.3}{\dB} & \SI{0.020}{\dB}    & no \\
    Ruido de disparo                         & \SI{-25.5}{\dB} & \SI{0.012}{\dB}    & no \\
    ASE del booster de transmisi\'on         & \SI{-37.9}{\dB} & \SI{0.001}{\dB}    & no \\
    \bottomrule
  \end{tabular}
  \caption[Contribuciones de ruido al receptor]{%
    Peso de cada fuente de ruido en el punto de operaci\'on. La segunda columna
    es la densidad espectral de la fuente respecto de la del batido
    se\~nal--ASE del pre-amplificador, que se toma como referencia. La tercera columna se refiere a
    cu\'anto subir\'ia el ruido total si se agregara ese t\'ermino al modelo.}
  \label{tab:ruidos}
\end{table}



\section{Desempe\~no te\'orico}
\label{sec:enl-teorico}

El umbral HD-FEC fija el requisito de OSNR. Utilizando las
conversiones desarrolladas en la \cref{sec:coh-metricas}, se obtiene la
\cref{tab:cadena-teorica}.

\begin{table}[htbp]
  \centering

  \begin{tabular}{llc}
    \toprule
    Magnitud & Origen & Valor \\
    \midrule
    $\BER$ objetivo & umbral HD-FEC       & \num{3.8e-3} \\
    $\EbNo$         & \cref{eq:ber-qpsk}  & \SI{5.52}{\dB} \\
    $\EsNo$         & $E_s = 2E_b$        & \SI{8.53}{\dB} \\
    $\OSNR$         & \cref{eq:osnr-esno} & \SI{12.61}{\dB} \\
    \bottomrule
  \end{tabular}
  \caption[Del umbral HD-FEC al requisito de OSNR]{%
    Conversiones desde el umbral de la correcci\'on de errores hasta
    el requisito de OSNR, a \SI{32}{\giga\baud} con
    $\Bref = \SI{12.5}{\giga\hertz}$.}
  \label{tab:cadena-teorica}
\end{table}

La misma cadena con una sola polarizaci\'on resulta en \SI{9.60}{\dB}, tres dB menos.
La diferencia sale de que multiplexar en polarizaci\'on obliga a duplicar la
potencia de se\~nal y deja la de ruido igual, as\'i que duplicar la tasa
binaria por esta v\'ia cuesta \SI{3.01}{\dB} de OSNR.

Los \SI{12.61}{\dB} son el desempe\~no de un receptor ideal. Lo que separa este
valor de OSNR con el que la simulaci\'on completa es
la penalidad de implementaci\'on, que se mide en la
\cref{sec:enl-verificacion}.
